% Compilar desde este directorio con:
% pdflatex Clase_04_Incertidumbre_en_agentes_inteligentes.tex
\documentclass[aspectratio=169,10pt]{beamer}

\usepackage[utf8]{inputenc}
\usepackage[T1]{fontenc}
\usepackage[spanish,es-nodecimaldot]{babel}
\usepackage{amsmath,amssymb,booktabs,tabularx,array,tikz}
\usetikzlibrary{arrows.meta,positioning,calc,fit,shapes.geometric}

\definecolor{UAPNavy}{HTML}{17324D}
\definecolor{UAPTeal}{HTML}{007F82}
\definecolor{UAPGold}{HTML}{E2A33A}
\definecolor{UAPLight}{HTML}{F3F6F8}
\definecolor{UAPGray}{HTML}{536471}
\definecolor{UAPRed}{HTML}{A94343}

\tikzset{
  flow/.style={-{Stealth[length=2.2mm]},thick,draw=UAPNavy},
  return/.style={-{Stealth[length=2mm]},thick,dashed,draw=UAPGray},
  datum/.style={draw=UAPGold,fill=UAPGold!16,rounded corners=2pt,align=center,minimum height=8mm,text=UAPNavy},
  process/.style={draw=UAPTeal,fill=UAPTeal!12,rounded corners=3pt,very thick,align=center,minimum height=10mm,text=UAPNavy},
  decision/.style={draw=UAPNavy,fill=UAPNavy!9,rounded corners=3pt,very thick,align=center,minimum height=10mm,text=UAPNavy},
  card/.style={draw=UAPTeal,fill=UAPLight,rounded corners=3pt,align=center,minimum height=11mm,text=UAPNavy},
  warning/.style={draw=UAPRed,fill=UAPRed!7,rounded corners=3pt,align=left},
  cell/.style={draw=UAPNavy,fill=white,minimum size=11mm,align=center,font=\bfseries,text=UAPNavy},
  hid/.style={cell,draw=UAPGold,fill=UAPGold!20,text=UAPNavy},
  go/.style={cell,draw=UAPTeal,fill=UAPTeal!15,text=UAPNavy}
}

\usetheme{Boadilla}
\usecolortheme{default}
\setbeamercolor{normal text}{fg=UAPNavy,bg=white}
\setbeamercolor{structure}{fg=UAPTeal}
\setbeamercolor{frametitle}{fg=white,bg=UAPNavy}
\setbeamercolor{title}{fg=white}
\setbeamercolor{subtitle}{fg=UAPGold}
\setbeamercolor{block title}{fg=white,bg=UAPTeal}
\setbeamercolor{block body}{fg=UAPNavy,bg=UAPLight}
\setbeamercolor{alertblock title}{fg=white,bg=UAPRed}
\setbeamercolor{alertblock body}{fg=UAPNavy,bg=UAPRed!7}
\setbeamerfont{frametitle}{series=\bfseries}
\setbeamertemplate{navigation symbols}{}
\setbeamertemplate{itemize item}{\color{UAPGold}$\blacktriangleright$}
\setbeamertemplate{itemize subitem}{\color{UAPTeal}$\bullet$}
\setbeamertemplate{footline}{%
  \leavevmode\hbox{%
    \begin{beamercolorbox}[wd=.78\paperwidth,ht=2.8ex,dp=1.2ex,leftskip=1em]{author in head/foot}%
      \color{white}\insertshorttitle
    \end{beamercolorbox}%
    \begin{beamercolorbox}[wd=.22\paperwidth,ht=2.8ex,dp=1.2ex,rightskip=1em plus1fil]{date in head/foot}%
      \color{UAPGray}\hfill Semana 4 \enspace | \enspace \insertframenumber/\inserttotalframenumber
    \end{beamercolorbox}%
  }%
}

\newcommand{\concepto}[1]{\textcolor{UAPTeal}{\textbf{#1}}}
\newcommand{\br}{\\}
\newcommand{\explicacion}[1]{\par\vspace{.25em}{\scriptsize\color{UAPGray}\textit{#1}\par}}
\newcommand{\pregunta}[1]{\begin{block}{Pregunta clave}#1\end{block}}

\title[IA · Clase 4]{Incertidumbre en agentes inteligentes}
\subtitle{De la observación imperfecta a la acción}
\author{Curso de Inteligencia Artificial}
\institute{Semana 4 · Clase 4}
\date{}

\begin{document}

% 1
\begin{frame}[plain]
  \begin{beamercolorbox}[wd=\paperwidth,ht=\paperheight,sep=2em]{frametitle}
    \vfill
    {\usebeamerfont{title}\usebeamercolor[fg]{title}\Huge\inserttitle\par}
    \vspace{.45em}
    {\usebeamerfont{subtitle}\usebeamercolor[fg]{subtitle}\Large\insertsubtitle\par}
    \vspace{1.1em}
    \begin{tikzpicture}[node distance=4mm,scale=.72,transform shape,every node/.style={text=UAPNavy}]
      \node[datum] (m) {Mundo\br real};
      \node[process,right=of m] (o) {Observación};
      \node[datum,right=of o] (c) {Creencia};
      \node[decision,right=of c] (d) {Acción};
      \node[datum,right=of d] (q) {Consecuencia};
      \draw[flow] (m)--(o); \draw[flow] (o)--(c); \draw[flow] (c)--(d); \draw[flow] (d)--(q);
      \draw[return] (q.south) to[out=-90,in=-90] node[below,font=\scriptsize,text=UAPGray]{actualización} (c.south);
    \end{tikzpicture}\par
    \vspace{.7em}
    {\color{white}\small El agente no conoce el mundo: lo percibe, forma una creencia, la actualiza y actúa asumiendo consecuencias.\par}
    \vfill
    {\color{UAPGold}\large\insertauthor\par}
    {\color{white}\insertinstitute\par}
  \end{beamercolorbox}
\end{frame}

% 2
\begin{frame}{Punto de partida}
  \small
  En las semanas 1 a 3 construimos agentes que resuelven problemas. La Semana 1 ya reconoció entornos \concepto{parcialmente observables} y \concepto{estocásticos}.
  \begin{itemize}
    \item Semana 1: agentes, entornos y racionalidad (observabilidad parcial, estocasticidad).
    \item Semana 2: representación de estados, acciones, objetivos y restricciones.
    \item Semana 3: búsqueda ciega e informada sobre un modelo fijo.
  \end{itemize}
  Para ejecutar esa búsqueda, las Semanas 2 y 3 adoptaron una \textbf{abstracción determinista}: el estado siempre era conocido con exactitud y cada acción llevaba a un único resultado posible.
  \begin{block}{Continuidad}
    Ahora levantamos esa abstracción: el agente \textbf{no observa todo}, sus sensores \textbf{fallan} y las acciones pueden tener \textbf{resultados múltiples}.
  \end{block}
\end{frame}

% 3
\begin{frame}{Propósito y resultados de aprendizaje}
  \small
  \textbf{Propósito:} comprender por qué aparece la incertidumbre, cómo se representa y cómo un agente decide sin fingir certeza.
  \begin{block}{Al finalizar, podremos}
    \begin{itemize}
      \item identificar fuentes y tipos de incertidumbre;
      \item distinguir estado real, percepción y creencia;
      \item construir una distribución conjunta y marginalizarla;
      \item actualizar una creencia con el teorema de Bayes;
      \item comparar acciones mediante costo esperado;
      \item reconocer cuándo una decisión se vuelve secuencial y qué papel juega una política;
      \item reconocer cómo se estiman estas probabilidades a partir de datos reales, más allá de una tabla dada.
    \end{itemize}
  \end{block}
  \explicacion{La clase es conceptual y numérica; no se implementan todavía algoritmos de aprendizaje por refuerzo.}
\end{frame}

% 4
\begin{frame}{Pregunta de apertura}
  \small
  \pregunta{Si un robot no puede saber con certeza si el camino está libre, ¿cómo decide cruzar?}
  \begin{columns}[T]
    \column{.60\textwidth}
      \begin{itemize}
        \item ¿Qué es exactamente lo que no sabe?
        \item ¿Qué puede observar y qué puede fallar?
        \item ¿Qué información tenía antes de observar?
        \item ¿Qué cuesta equivocarse en cada sentido?
        \item ¿Cómo cambia su decisión al recibir una señal?
      \end{itemize}
    \column{.36\textwidth}
      \centering
      \begin{tikzpicture}[node distance=2mm,scale=.78,transform shape]
        \node[cell] (s) {Inicio};
        \node[hid,right=of s] (m) {?};
        \node[go,right=of m] (g) {Meta};
        \draw[flow] (s)--(m); \draw[flow] (m)--(g);
      \end{tikzpicture}\par
      {\scriptsize\color{UAPGray} nuestro caso conductor, desarrollado en detalle más adelante}
  \end{columns}
  \explicacion{La incertidumbre no es un defecto del agente: es una propiedad de su relación con el entorno.}
\end{frame}

% 5
\begin{frame}{¿Por qué aparece la incertidumbre?}
  \small
  Un agente actúa sobre un mundo que no puede conocer por completo. Entre el mundo y la decisión median percepciones, modelos y evidencia, cada uno con pérdidas.
  \vspace{.4em}
  \centering
  \begin{tikzpicture}[node distance=5mm,scale=.78,transform shape]
    \node[datum] (m) {Mundo\br real};
    \node[process,right=of m] (o) {Observación};
    \node[datum,right=of o] (c) {Creencia};
    \node[decision,right=of c] (a) {Acción};
    \node[datum,right=of a] (q) {Consecuencia};
    \draw[flow] (m)--(o); \draw[flow] (o)--(c); \draw[flow] (c)--(a); \draw[flow] (a)--(q);
    \node[warning,below=9mm of c,text width=64mm,font=\scriptsize] {\textbf{Pérdidas}: lo no observado, el ruido, lo mal modelado y lo que cambiará en el futuro.};
  \end{tikzpicture}
  \vspace{.55em}
  \begin{alertblock}{Idea}
    Decidir con certeza es un caso especial; lo habitual es decidir con información parcial.
  \end{alertblock}
\end{frame}

% 6
\begin{frame}{Fuentes de incertidumbre}
  \small
  \begin{columns}[T]
    \column{.48\textwidth}
      \begin{block}{Del lado de la observación}
        \begin{itemize}
          \item \textbf{Observabilidad parcial}: el agente no ve todo el estado.
          \item \textbf{Ruido de sensores}: la lectura no es el estado.
          \item \textbf{Información faltante o ambigua}.
          \item \textbf{Datos insuficientes o mal cubiertos}.
        \end{itemize}
      \end{block}
    \column{.48\textwidth}
      \begin{block}{Del lado del mundo y del modelo}
        \begin{itemize}
          \item \textbf{Resultados estocásticos}: una acción, varios efectos posibles.
          \item \textbf{Estados omitidos}: el modelo simplifica la realidad.
          \item \textbf{Cambio del entorno}: el futuro difiere del pasado.
        \end{itemize}
      \end{block}
  \end{columns}
  \vfill
  \explicacion{Las fuentes se combinan: un sensor ruidoso sobre un mundo cambiante y modelado de forma incompleta.}
\end{frame}

% 7
\begin{frame}{Dos tipos fundamentales}
  \small
  \begin{columns}[T]
    \column{.48\textwidth}
      \begin{block}{Incertidumbre aleatoria}
        Variabilidad que \textbf{permanece} aun conociendo el modelo exacto.
        \begin{itemize}
          \item El resultado de un dado o de un sensor ruidoso.
          \item El resultado de una acción con efectos múltiples.
        \end{itemize}
        \textbf{Respuesta}: reservas, redundancia, políticas robustas.
      \end{block}
    \column{.48\textwidth}
      \begin{block}{Incertidumbre epistémica}
        Falta de conocimiento \textbf{potencialmente reducible}.
        \begin{itemize}
          \item Pocos datos o parámetros imprecisos.
          \item Modelo incompleto o mal especificado.
        \end{itemize}
        \textbf{Respuesta}: más datos pertinentes, mejores mediciones.
      \end{block}
  \end{columns}
  \vfill
  \alert{La frontera depende del modelo: lo que parece aleatorio puede explicarse al añadir variables.}
\end{frame}

% 8
\begin{frame}{Otras distinciones útiles}
  \scriptsize
  \begin{tabularx}{\textwidth}{>{\bfseries}p{.26\textwidth}X}
    \toprule
    Tipo & Pregunta que responde y ejemplo \\\midrule
    Observacional & ¿La medición coincide con el estado? El sensor reporta \texttt{alerta} aunque el camino esté libre. \\
    Estructural & ¿El modelo incluye lo relevante? Un robot que ignora que el piso puede estar húmedo. \\
    Distribucional\br (cambio de distribución) & ¿El futuro se parece al pasado? Condiciones de hoy distintas de las históricas. \\
    Vaguedad & ¿La categoría tiene límites nítidos? ``Cerca'', ``alto'' o ``riesgoso'' no son lo mismo que aleatorio. \\
    \bottomrule
  \end{tabularx}
  \vfill
  \explicacion{Ninguna etiqueta es universal: clasificar sirve para elegir la respuesta, no para fijar una división metafísica.}
\end{frame}

% 9
\begin{frame}{Problemáticas asociadas (I)}
  \small
  La incertidumbre no resuelta produce errores de razonamiento:
  \begin{itemize}
    \item \concepto{Confundir ``no sé'' con ``no existe''}: la ausencia de dato no prueba ausencia del fenómeno.
    \item \concepto{Invertir la condición}: $P(\text{alerta}\mid \text{bloqueo}) \neq P(\text{bloqueo}\mid \text{alerta})$.
    \item \concepto{Ignorar la tasa base}: una señal poco común puede seguir indicando un estado poco probable.
    \item \concepto{Contar dos veces evidencia dependiente}: dos sensores con una causa común no aportan información independiente.
  \end{itemize}
  \vfill
  \explicacion{Cada error convierte un número correcto en una conclusión equivocada.}
\end{frame}

% 10
\begin{frame}{Problemáticas asociadas (II)}
  \small
  \begin{itemize}
    \item \concepto{Falsa precisión}: asignar un número sin respaldo ni análisis de sensibilidad.
    \item \concepto{Confundir probabilidad con importancia}: un evento raro puede ser gravísimo.
    \item \concepto{Ocultar incertidumbre del modelo}: presentar una creencia como si fuera un hecho.
    \item \concepto{Actuar fuera de distribución}: confiar en datos históricos ante condiciones nuevas.
    \item \concepto{Explorar sin control}: aprender por prueba y error en un entorno donde equivocarse daña.
  \end{itemize}
  \vfill
  \explicacion{El riesgo no está en usar probabilidades, sino en presentarlas como certezas que no son.}
\end{frame}

% 11
\begin{frame}{¿Qué podemos hacer?}
  \small
  \begin{columns}[T]
    \column{.48\textwidth}
      \begin{block}{Reducir y representar}
        \begin{itemize}
          \item Obtener información pertinente.
          \item Modelar estados posibles y sus probabilidades.
          \item Mantener conjuntos de estados consistentes.
        \end{itemize}
      \end{block}
      \begin{block}{Actualizar}
        \begin{itemize}
          \item Incorporar evidencia con Bayes.
          \item Combinar observaciones sin doble conteo.
        \end{itemize}
      \end{block}
    \column{.48\textwidth}
      \begin{block}{Decidir y controlar}
        \begin{itemize}
          \item Comparar acciones por costo o utilidad esperada.
          \item Hacer análisis de sensibilidad.
          \item Elegir políticas robustas ante escenarios.
        \end{itemize}
      \end{block}
      \begin{block}{Limitar el daño}
        \begin{itemize}
          \item Abstenerse o pedir intervención humana.
          \item Aprender por interacción solo si es seguro.
        \end{itemize}
      \end{block}
  \end{columns}
\end{frame}

% 12
\begin{frame}{El flujo completo del agente}
  \small
  \centering
  \begin{tikzpicture}[node distance=5mm,scale=.74,transform shape]
    \node[datum] (m) {Mundo\br real};
    \node[process,right=of m] (o) {Observación};
    \node[datum,right=of o] (b) {Creencia\br $P(\text{estado}\mid E)$};
    \node[decision,right=of b] (a) {Acción};
    \node[datum,right=of a] (q) {Consecuencia\br y costo};
    \draw[flow] (m)--(o); \draw[flow] (o)--(b); \draw[flow] (b)--(a); \draw[flow] (a)--(q);
    \draw[return] (q.south) to[out=-90,in=-90] (b.south);
  \end{tikzpicture}
  \vfill
  \begin{itemize}
    \item La \concepto{creencia} resume todo lo que el agente sabe del estado.
    \item La \concepto{acción} se elige comparando consecuencias bajo esa creencia.
    \item La \concepto{actualización} convierte la observación nueva en creencia nueva.
  \end{itemize}
  \explicacion{Esta cadena organiza todo lo que sigue: representación, actualización y decisión.}
\end{frame}

% 13
\begin{frame}{Caso conductor: un robot con sensor}
  \small
  Un robot avanza por un pasillo de tres celdas y debe llegar a la meta. La celda central puede estar \textbf{libre} o \textbf{bloqueada}; el robot no lo sabe.
  \vspace{.45em}
  \centering
  \begin{tikzpicture}[node distance=3mm]
    \node[cell] (s) {Inicio};
    \node[hid,right=of s] (m) {?};
    \node[go,right=of m] (g) {Meta};
    \draw[flow] (s)--(m); \draw[flow] (m)--(g);
    \node[font=\scriptsize,text=UAPGray,above=4mm of m] {estado oculto};
  \end{tikzpicture}
  \vfill
  \begin{itemize}
    \item Un sensor informa \texttt{alerta} o \texttt{sin alerta} sobre la celda central.
    \item El sensor es \textbf{imperfecto}: a veces se equivoca.
    \item El robot elige entre \textbf{cruzar} directo o \textbf{rodear} por un camino más largo.
  \end{itemize}
  \explicacion{Este escenario mínimo aísla los conceptos: estado oculto, evidencia ruidosa y decisión bajo costo.}
\end{frame}

% 14
\begin{frame}{¿Qué desconoce el robot?}
  \small
  El robot conoce su posición y su meta, pero ignora el \concepto{estado real} de la celda central.
  \begin{block}{Separación fundamental}
    \begin{itemize}
      \item \textbf{Estado real}: libre o bloqueada, exista o no quien lo observe.
      \item \textbf{Percepción}: la lectura del sensor (\texttt{alerta} o \texttt{sin alerta}).
      \item \textbf{Creencia}: lo que el robot concluye a partir de sus observaciones previas.
    \end{itemize}
  \end{block}
  \alert{El estado no cambia por no observarlo; cambia lo que el agente puede afirmar sobre él.}
\end{frame}

% 15
\begin{frame}{Estado real vs percepción}
  \small
  La lectura del sensor no es el estado: es una \concepto{proyección ruidosa}.
  \begin{center}
  \begin{tabular}{ll}
    \toprule
    $B$ & evento ``la celda está bloqueada'' \\
    $L$ & evento ``la celda está libre'' \\
    $A$ & lectura ``el sensor informa alerta'' \\
    $\neg A$ & lectura ``el sensor informa sin alerta'' \\
    \bottomrule
  \end{tabular}
  \end{center}
  \explicacion{La dirección de la condición cambia la pregunta: $P(A\mid B)$ y $P(A\mid L)$ parten del estado real y describen qué lectura produce el sensor; $P(B\mid A)$ parte de una alerta observada y expresa cuánta probabilidad merece el bloqueo. En $P(A\mid L)$, $A$ sigue significando ``el sensor emite una alerta'': como el estado real es libre, esa alerta es un falso positivo.}
\end{frame}

% 16
\begin{frame}{El sensor y la previa}
  \small
  Para cuantificar el caso necesitamos dos ingredientes: cómo se comporta el sensor y qué tan probable era el bloqueo antes de mirar. ``Alerta'' es la salida del sensor, no el estado.
  \begin{columns}[T]
    \column{.48\textwidth}
      \begin{block}{El sensor (dado el estado)}
        \[
          P(A\mid B)=0{,}90 \qquad P(A\mid L)=0{,}10
        \]
        \begin{itemize}
          \item \textbf{Sensibilidad} $P(A\mid B)$: si está bloqueada, el sensor alerta en 90 de cada 100 casos.
          \item \textbf{Falsos positivos} $P(A\mid L)$: si está libre, alerta igual en 10 de cada 100; son alertas falsas.
        \end{itemize}
      \end{block}
    \column{.48\textwidth}
      \begin{block}{La previa (antes de mirar)}
        \[
          P(B)=0{,}20 \qquad P(L)=0{,}80
        \]
        \begin{itemize}
          \item Antes de observar, 20 de cada 100 celdas están bloqueadas.
          \item Esta frecuencia previa condiciona cuánto cambia la creencia.
        \end{itemize}
      \end{block}
  \end{columns}
  \begin{alertblock}{Valores ilustrativos}
    $P(B)=0{,}20$, $P(A\mid B)=0{,}90$ y $P(A\mid L)=0{,}10$ se eligen para el ejercicio; en un caso real se estiman con datos o especificaciones del sensor.
  \end{alertblock}
\end{frame}

% 17
\begin{frame}{Cómo se construye la tabla}
  \small
  Cada celda de la tabla combina la \concepto{previa} con el comportamiento del sensor:
  \[
    P(\text{estado},\text{lectura})=P(\text{estado})\,P(\text{lectura}\mid \text{estado})
  \]
  Primero, los complementos del sensor:
  \[
    P(\neg A\mid B)=1-0{,}90=0{,}10, \qquad P(\neg A\mid L)=1-0{,}10=0{,}90
  \]
  \begin{columns}[T]
    \column{.48\textwidth}
      \[
        P(B,A)=P(B)\,P(A\mid B)=0{,}20(0{,}90)=0{,}18
      \]
      \[
        P(L,A)=P(L)\,P(A\mid L)=0{,}80(0{,}10)=0{,}08
      \]
    \column{.48\textwidth}
      \[
        P(B,\neg A)=P(B)\,P(\neg A\mid B)=0{,}20(0{,}10)=0{,}02
      \]
      \[
        P(L,\neg A)=P(L)\,P(\neg A\mid L)=0{,}80(0{,}90)=0{,}72
      \]
  \end{columns}
  \explicacion{Cada expresión es una celda de la tabla siguiente: $P(B,A)$ bloqueada \emph{con} alerta (verdadero positivo), $P(B,\neg A)$ bloqueada \emph{sin} alerta (falso negativo), $P(L,A)$ libre \emph{con} alerta (falso positivo) y $P(L,\neg A)$ libre \emph{sin} alerta (verdadero negativo).}
\end{frame}

% 18
\begin{frame}{Cómo se interpreta la tabla}
  \small
  \begin{center}
  \begin{tabular}{lrrr}
    \toprule
    & $A$ (alerta) & $\neg A$ (sin alerta) & Marginal \\\midrule
    $B$ (bloqueada) & $0{,}18$ & $0{,}02$ & $0{,}20$ \\
    $L$ (libre) & $0{,}08$ & $0{,}72$ & $0{,}80$ \\
    \midrule
    Marginal & $0{,}26$ & $0{,}74$ & $1{,}00$ \\
    \bottomrule
  \end{tabular}
  \end{center}
  \begin{itemize}
    \item \concepto{Filas}: estado real ($B$ bloqueada, $L$ libre).
    \item \concepto{Columnas interiores}: lectura del sensor ($A$ alerta, $\neg A$ sin alerta).
    \item \concepto{Celda interior}: probabilidad conjunta $P(\text{estado},\text{lectura})$.
    \item \concepto{Margen de fila}: previas $P(B)$ y $P(L)$.
    \item \concepto{Margen de columna}: $P(A)$ y $P(\neg A)$, la probabilidad de cada lectura.
    \item \concepto{Esquina}: suma total $1$, comprobación de normalización.
  \end{itemize}
  \explicacion{$P(B,A)=0{,}18$: probabilidad de que la celda esté bloqueada \emph{y} el sensor informe alerta. Que exista una alerta no demuestra bloqueo: la celda $P(L,A)=0{,}08$ muestra que el sensor también alerta cuando la celda está libre.}
\end{frame}

% 19
\begin{frame}{Marginalizar: probabilidad de una lectura}
  \small
  La probabilidad de una lectura se obtiene \concepto{marginalizando} el estado:
  \[
    P(A)=P(A\mid B)P(B)+P(A\mid L)P(L)
  \]
  Con los números del ejemplo:
  \[
    P(A)=0{,}90(0{,}20)+0{,}10(0{,}80)=0{,}18+0{,}08=0{,}26
  \]
  \[
    P(\neg A)=0{,}10(0{,}20)+0{,}90(0{,}80)=0{,}02+0{,}72=0{,}74
  \]
  \explicacion{Para saber cuán probable es una alerta se suman todos los casos en que el sensor alerta: bloqueada con alerta más libre con alerta, $P(A)=P(B,A)+P(L,A)=0{,}18+0{,}08=0{,}26$.}
\end{frame}

% 20
\begin{frame}{Independencia de observaciones}
  \small
  Dos observaciones son \concepto{independientes} (dado el estado) cuando cada una aporta información \textbf{nueva}:
  \[
    P(A_1,A_2\mid B)=P(A_1\mid B)\,P(A_2\mid B)
  \]
  \begin{itemize}
    \item Si dos sensores comparten una \textbf{causa común} o miden el \textbf{mismo fallo}, multiplicarlos cuenta dos veces la evidencia.
    \item Asumir independencia cuando no existe produce \concepto{confianza excesiva}.
  \end{itemize}
  \begin{block}{En nuestro caso}
    Si el robot incorporara un segundo sensor, esta sería la condición para combinar ambas lecturas sin sobreestimar la evidencia.
  \end{block}
  \explicacion{La independencia es un supuesto que debe justificarse, no una propiedad automática de tener dos mediciones.}
\end{frame}

% 21
\begin{frame}{Actualizar creencias: teorema de Bayes}
  \small
  La creencia sobre el estado, tras observar la lectura, se actualiza:
  \[
    P(B\mid A)=\frac{P(A\mid B)\,P(B)}{P(A)}
  \]
  \begin{itemize}
    \item $P(B)$ \concepto{(previa)}: ¿qué probabilidad tenía el bloqueo antes de observar?
    \item $P(A\mid B)$ \concepto{(verosimilitud)}: si había bloqueo, ¿qué probabilidad tenía el sensor de alertar?
    \item $P(A)$ \concepto{(evidencia)}: ¿qué probabilidad había de observar una alerta, cualquiera fuera el estado?
    \item $P(B\mid A)$ \concepto{(posterior)}: después de observar la alerta, ¿qué probabilidad tiene el bloqueo?
  \end{itemize}
  \explicacion{Bayes no decide por nosotros: convierte la observación en una creencia coherente con el modelo.}
\end{frame}

% 22
\begin{frame}{Efecto de la tasa base}
  \small
  Con una previa $P(B)=0{,}20$:
  \[
    P(B\mid A)=\frac{0{,}90(0{,}20)}{0{,}26}\approx0{,}6923
  \]
  \begin{block}{Observación crítica}
    Aunque la sensibilidad es alta, la posterior no es $0{,}90$: una alerta puede venir de una celda bloqueada o de una libre, y la frecuencia previa de cada estado decide cuántas alertas corresponden a bloqueos reales.
  \end{block}
  Si la celda bloqueada fuera \textbf{rarísima}, por ejemplo $P(B)=0{,}01$, una alerta apenas elevaría la creencia:
  \[
    P(B\mid A)=\frac{0{,}90(0{,}01)}{0{,}90(0{,}01)+0{,}10(0{,}99)}\approx0{,}083
  \]
  \explicacion{La misma señal no significa lo mismo en poblaciones con prevalencias distintas.}
\end{frame}

% 23
\begin{frame}{Números del ejemplo}
  \small
  \begin{center}
  \begin{tabular}{lll}
    \toprule
    Lectura & Posterior $P(B\mid \cdot)$ & Lectura de la creencia \\\midrule
    $\neg A$ (sin alerta) & $\dfrac{0{,}10(0{,}20)}{0{,}74}\approx0{,}027$ & creencia baja en bloqueo \\
    $A$ (alerta) & $\dfrac{0{,}90(0{,}20)}{0{,}26}\approx0{,}692$ & creencia alta en bloqueo \\
    \bottomrule
  \end{tabular}
  \end{center}
  \vfill
  \begin{itemize}
    \item Sin alerta, el bloqueo pasa de $0{,}20$ a $0{,}027$.
    \item Con alerta, el bloqueo pasa de $0{,}20$ a $0{,}692$.
  \end{itemize}
  \explicacion{La misma previa produce posteriores opuestas según la observación; ahora falta convertir creencia en acción.}
\end{frame}

% 24
\begin{frame}{De la creencia a la acción: matriz de costos}
  \small
  La creencia por sí sola no dice qué hacer; hace falta valorar las \concepto{consecuencias}.
  \begin{center}
  \begin{tabular}{lrr}
    \toprule
    Acción & Bloqueada & Libre \\\midrule
    Cruzar & 30 & 0 \\
    Rodear & 8 & 8 \\
    \bottomrule
  \end{tabular}
  \end{center}
  \begin{itemize}
    \item Cruzar es gratis si la celda está libre, pero cuesta 30 si choca.
    \item Rodear siempre cuesta 8: es el desvío seguro.
  \end{itemize}
  \explicacion{Valores hipotéticos del ejercicio; en un caso real se fijan según daño, tiempo y recursos. Un falso negativo (cruzar con bloqueo) puede costar más que un falso positivo (rodear sin bloqueo); cada error se valora por separado.}
\end{frame}

% 25
\begin{frame}{Valor esperado de una acción}
  \small
  Usamos $EC$ por \textit{Expected Cost} (\concepto{costo esperado}). Para una acción $a$, evidencia $E$ y posibles estados $s$:
  \[
    EC(a\mid E)=\sum_{s}P(s\mid E)\,C(a,s)
  \]
  donde $C(a,s)$ es el costo de ejecutar $a$ cuando ocurre $s$.
  \br
  En criollo: para cada estado posible $s$ multiplicamos qué tan probable creemos que es ($P(s\mid E)$) por lo que costaría actuar si ese estado fuera el real ($C(a,s)$), y sumamos. El resultado es un promedio de costos ponderado por la creencia: no predice qué va a pasar, mide qué tan cara es la acción \textbf{en promedio}.
  \explicacion{La evidencia $E$ puede ser una lectura del sensor, pero también puede no haber evidencia todavía. La fórmula es siempre la misma; lo que cambia de un caso a otro es qué creencia $P(s\mid E)$ usamos.}
\end{frame}

% 26
\begin{frame}{Decidir sin haber observado nada}
  \small
  Antes de leer el sensor no hay evidencia: $E=\varnothing$. En ese caso $P(s\mid E)$ es simplemente la previa $P(s)$, porque todavía no condicionamos en ningún dato.
  \br
  Con previa $P(B)=0{,}20$ (y $P(L)=0{,}80$), reemplazando en la fórmula:
  \[
    EC(\text{cruzar})=P(B)(30)+P(L)(0)=0{,}20(30)+0{,}80(0)=6
  \]
  \[
    EC(\text{rodear})=P(B)(8)+P(L)(8)=8
  \]
  \begin{block}{Resultado}
    Con la previa, \textbf{cruzar} tiene menor costo esperado ($6<8$).
  \end{block}
  \explicacion{Es el mismo $EC(a\mid E)$ de la diapositiva anterior, solo que con $E=\varnothing$: por eso ``sin observar'' no usa otra fórmula, usa la misma con la creencia inicial.}
\end{frame}

% 27
\begin{frame}{Decidir después de observar}
  \small
  Con alerta, la creencia sube a $P(B\mid A)=0{,}692$:
  \[
    EC(\text{cruzar}\mid A)=0{,}692(30)+0{,}308(0)\approx20{,}8
  \]
  \[
    EC(\text{rodear}\mid A)=8
  \]
  Con sin alerta, la creencia baja a $P(B\mid \neg A)=0{,}027$:
  \[
    EC(\text{cruzar}\mid \neg A)=0{,}027(30)+0{,}973(0)\approx0{,}8
  \]
  \begin{block}{Regla resultante}
    \textbf{Rodear} ante alerta; \textbf{cruzar} ante sin alerta.
  \end{block}
  \explicacion{La observación cambió la acción: ahora la decisión es una función de la evidencia, no una elección fija.}
\end{frame}

% 28
\begin{frame}{Umbral y sensibilidad}
  \small
  Llamemos $p=P(B\mid E)$ a la creencia de bloqueo, con o sin evidencia. Cruzar conviene cuando su costo esperado es menor que rodear:
  \[
    p(30)+(1-p)(0)<8 \quad\Longrightarrow\quad p<\frac{8}{30}\approx0{,}267
  \]
  \begin{itemize}
    \item Previa $p=0{,}20<0{,}267$: conviene cruzar sin observar.
    \item Alerta $p=0{,}692>0{,}267$: conviene rodear.
    \item Sin alerta $p=0{,}027<0{,}267$: conviene cruzar.
  \end{itemize}
  \begin{alertblock}{Análisis de sensibilidad}
    Si los costos cambian (choque más grave, desvío más barato), el umbral se desplaza y la decisión puede invertirse.
  \end{alertblock}
\end{frame}

% 29
\begin{frame}{Errores frecuentes al decidir}
  \small
  \begin{itemize}
    \item \concepto{Invertir el condicional}: tratar $P(A\mid B)$ como si fuera $P(B\mid A)$.
    \item \concepto{Ignorar la previa}: usar solo la sensibilidad del sensor.
    \item \concepto{Confundir creencia con certeza}: ``el sensor dijo alerta'' no es ``está bloqueada''.
    \item \concepto{Valorar mal los errores}: tratar falso positivo y falso negativo como equivalentes.
    \item \concepto{Omitir el costo de actuar}: decidir sin sumar el precio de la propia decisión.
  \end{itemize}
  \explicacion{La fórmula es corta; construir bien sus componentes es el trabajo real.}
\end{frame}

% 30
\begin{frame}{De una decisión a una secuencia}
  \small
  Hasta aquí decidimos una sola vez. En un problema real, cada acción \concepto{cambia el estado} y condiciona las decisiones siguientes.
  \vspace{.4em}
  \begin{itemize}
    \item Cruzar coloca al robot en otra celda; rodear consume batería.
    \item La posición, la batería y el riesgo acumulado forman el \concepto{estado}.
    \item Ya no basta ``la mejor acción ahora'': importa la mejor \concepto{regla} para cada situación futura.
  \end{itemize}
  \begin{block}{Nuevo objeto}
    Una \textbf{política} indica qué acción tomar en cada estado, anticipando consecuencias futuras.
  \end{block}
  \explicacion{La decisión única es el caso especial en que el futuro no depende de la acción presente.}
\end{frame}

% 31
\begin{frame}{Política: qué hacer en cada situación}
  \small
  Con observabilidad parcial, el agente no conoce el estado oculto: \concepto{decide sobre su creencia} acerca de él, junto con lo que sí conoce (posición, batería).
  \begin{center}
  \begin{tabular}{ll}
    \toprule
    Situación (posición, batería, creencia) & Acción de la política \\\midrule
    creencia de bloqueo por debajo del umbral ($<0{,}267$) & cruzar \\
    creencia por encima del umbral, batería suficiente & rodear \\
    creencia alta, batería insuficiente para rodear & detenerse / pedir asistencia \\
    \bottomrule
  \end{tabular}
  \end{center}
  \begin{itemize}
    \item La política condensa la decisión en una regla por situación.
    \item Puede ser determinista (siempre la misma acción) o estocástica (con probabilidades).
  \end{itemize}
  \explicacion{Una política no es indecisión ni una lista de ocurrencias: es una regla sistemática y auditable.}
\end{frame}

% 32
\begin{frame}{Recompensa y exploración}
  \small
  Si el agente debe aprender a actuar por interacción, aparece el \concepto{aprendizaje por refuerzo} (RL).
  \begin{block}{Dos tensiones}
    \begin{itemize}
      \item \textbf{Recompensa}: expresa el objetivo, pero puede estar mal diseñada y premiar atajos.
      \item \textbf{Explorar vs explotar}: probar acciones da información y también costo o riesgo.
    \end{itemize}
  \end{block}
  \begin{itemize}
    \item Explorar una celda desconocida puede revelar un obstáculo o provocar un choque.
    \item Explotar siempre lo conocido consolida una creencia quizá equivocada.
  \end{itemize}
\end{frame}

% 33
\begin{frame}{¿Cuándo corresponde RL?}
  \small
  RL no es la respuesta a toda incertidumbre. Es apropiado cuando:
  \begin{itemize}
    \item hay decisiones \concepto{secuenciales} significativas;
    \item la acción presente \concepto{modifica} estados y opciones futuras;
    \item existe \concepto{retroalimentación} repetida y un objetivo medible;
    \item se puede \concepto{explorar con seguridad} o en simulación.
  \end{itemize}
  \vfill
  \begin{alertblock}{No usar RL cuando}
    \begin{itemize}
      \item una regla simple ya resuelve el problema;
      \item la interacción es escasa, costosa o no ética;
      \item equivocarse durante el entrenamiento es inaceptable.
    \end{itemize}
  \end{alertblock}
\end{frame}

% 34
\begin{frame}{Transferencia: el caso de una app de rutas}
  \small
  Una app de rutas debe elegir entre un camino \textbf{largo}, siempre disponible, y un camino \textbf{corto} cuyo estado no puede confirmar con certeza: se nutre de las trazas GPS de los celulares que lo cruzan, y por allí pasa tan poca gente que la señal es débil.

  \vspace{.3em}
  \centering
  \begin{tikzpicture}[scale=.85,transform shape]
    \node[cell] (s) {Origen};
    \node[go,right=50mm of s] (g) {Destino};
    \draw[flow] (s) to[bend left=25] node[above,font=\scriptsize,align=center,text width=40mm] {camino corto: ¿transitable u obstruido?} (g);
    \draw[flow] (s) to[bend right=25] node[below,font=\scriptsize] {camino largo: siempre disponible} (g);
  \end{tikzpicture}
  \vfill
  \begin{itemize}
    \item \textbf{Estado real} $\leftrightarrow$ el paso corto es transitable en auto (aunque sea lento) u obstruido para autos.
    \item \textbf{Lectura} $\leftrightarrow$ trazas GPS de otros vehículos en ese tramo, en los últimos minutos.
    \item \textbf{Previa} $\leftrightarrow$ tasa histórica de obstrucción del paso, para esa franja horaria.
    \item \textbf{Matriz de costos} $\leftrightarrow$ tiempo de viaje esperado, en minutos.
  \end{itemize}
  \explicacion{El mismo esquema del robot ---estado, lectura, previa, sensor y costo--- se aplica; lo que cambia es el dominio.}
\end{frame}

% 35
\begin{frame}{El sensor y la previa en el cruce}
  \small
  Igual que con el robot, hacen falta dos ingredientes: cómo se comporta la evidencia y qué tan probable era la obstrucción antes de mirar.
  \begin{columns}[T]
    \column{.48\textwidth}
      \begin{block}{La evidencia GPS (dado el estado)}
        \[
          P(D\mid O)=0{,}95 \qquad P(D\mid T)=0{,}30
        \]
        \begin{itemize}
          \item $D$: la señal es débil (pocas o ninguna traza vehicular reciente en el tramo).
          \item $P(D\mid T)=0{,}30$: si el camino está transitable, en 3 de cada 10 ventanas la señal es igualmente débil; es una falsa señal de obstrucción.
        \end{itemize}
      \end{block}
    \column{.48\textwidth}
      \begin{block}{La previa (antes de mirar)}
        \[
          P(O)=0{,}30 \qquad P(T)=0{,}70
        \]
        \begin{itemize}
          \item Tasa histórica de obstrucción del paso corto en esa franja horaria.
          \item $O$: obstruido para autos. $T$: transitable, aunque sea lento.
        \end{itemize}
      \end{block}
  \end{columns}
  \begin{alertblock}{Diferencia con el robot}
    Acá la evidencia es menos discriminante: $P(D\mid T)=0{,}30$ no es chico, porque el bajo tránsito propio de la calle también genera poca señal. Valores ilustrativos; en un caso real se estiman con datos históricos del propio sistema.
  \end{alertblock}
\end{frame}

% 36
\begin{frame}{Tabla conjunta y posterior}
  \small
  Con el mismo procedimiento de las diapositivas 17 a 21 ($P(\text{estado},\text{lectura})=P(\text{estado})\,P(\text{lectura}\mid\text{estado})$, y Bayes para actualizar la probabilidad de obstrucción):
  \begin{center}
  \begin{tabular}{lrrr}
    \toprule
    & $D$ (señal débil) & $\neg D$ (señal normal) & Marginal \\\midrule
    $O$ (obstruido) & $0{,}285$ & $0{,}015$ & $0{,}30$ \\
    $T$ (transitable) & $0{,}21$ & $0{,}49$ & $0{,}70$ \\
    \midrule
    Marginal & $0{,}495$ & $0{,}505$ & $1{,}00$ \\
    \bottomrule
  \end{tabular}
  \end{center}
  \[
    P(O\mid D)=\frac{0{,}285}{0{,}495}\approx0{,}576 \qquad\qquad P(O\mid \neg D)=\frac{0{,}015}{0{,}505}\approx0{,}030
  \]
  \explicacion{Con señal débil, la creencia de obstrucción sube de $0{,}30$ a $0{,}576$; con señal normal, baja a $0{,}030$.}
\end{frame}

% 37
\begin{frame}{Costos y decisión: ¿corto o largo?}
  \small
  \begin{center}
  \begin{tabular}{lrr}
    \toprule
    Ruta & Transitable & Obstruido \\\midrule
    Corto & 5 & 20 \\
    Largo & 12 & 12 \\
    \bottomrule
  \end{tabular}
  \end{center}
  {\scriptsize\color{UAPGray} minutos de viaje; si el corto resulta obstruido, hay que retroceder y tomar igual el largo.}
  \[
    EC(\text{corto}\mid D)=5(0{,}424)+20(0{,}576)\approx13{,}6 \qquad EC(\text{largo}\mid D)=12
  \]
  \[
    EC(\text{corto}\mid \neg D)=5(0{,}970)+20(0{,}030)\approx5{,}4 \qquad EC(\text{largo}\mid \neg D)=12
  \]
  \begin{block}{Regla resultante}
    \textbf{Largo} ante señal débil ($13{,}6>12$); \textbf{corto} ante señal normal ($5{,}4<12$) --- la misma estructura que ``rodear ante alerta, cruzar ante sin alerta''.
  \end{block}
  \explicacion{El umbral de indiferencia es $p^{*}=7/15\approx0{,}467$: exactamente donde cae la frontera entre ambas reglas.}
\end{frame}

% 38
\begin{frame}{De la tabla a un sistema real}
  \small
  En el robot, la tabla venía dada. En un sistema real hay que construirla:
  \begin{itemize}
    \item \textbf{Previa}: se estima con el historial del propio sistema, condicionada a franja horaria y día --- agregarla sin condicionar repite el error de ignorar la tasa base.
    \item \textbf{Evidencia dado el estado}: se estima contando frecuencias sobre pares (señal observada, estado confirmado después), no se define a mano.
    \item \textbf{Pocos datos, sensor débil}: en un tramo de bajo tránsito hay pocas observaciones para calibrar; esa es incertidumbre epistémica sobre el propio modelo, no solo sobre el estado.
    \item \textbf{En producción}: no se arma una tabla de dos estados; se entrena un clasificador que combina varias señales y devuelve $P(\text{obstruido}\mid\text{features})$ directamente --- el mismo cálculo, aprendido de datos en vez de definido a mano.
    \item \textbf{Combinar evidencia y actualizar}: sumar fuentes exige justificar independencia (diapositiva 20); si el patrón de uso cambia, el modelo queda desactualizado (actuar fuera de distribución, diapositiva 10).
  \end{itemize}
  \explicacion{La mecánica de Bayes no cambia; lo que cambia es de dónde provienen los parámetros probabilísticos que utiliza.}
\end{frame}

% 39
\begin{frame}{También aplica: el caso del agua}
  \small
  El mismo patrón ---estado real, lectura, previa, evidencia y matriz de costos--- se aplica a la vigilancia de calidad del agua: el estado es la condición real del agua, la lectura es el resultado del análisis o del sensor, y la decisión es revisar o no revisar.
\end{frame}

% 40
\begin{frame}{Síntesis conceptual}
  \small
  \begin{center}
  \begin{tikzpicture}[node distance=5mm,scale=.72,transform shape]
    \node[datum] (m) {Mundo\br real};
    \node[process,right=of m] (o) {Observación};
    \node[datum,right=of o] (b) {Creencia};
    \node[decision,right=of b] (a) {Acción};
    \node[datum,right=of a] (q) {Consecuencia};
    \draw[flow] (m)--(o); \draw[flow] (o)--(b); \draw[flow] (b)--(a); \draw[flow] (a)--(q);
    \draw[return] (q.south) to[out=-90,in=-90] (b.south);
  \end{tikzpicture}
  \end{center}
  \vfill
  \begin{itemize}
    \item La incertidumbre se \concepto{origina} en observación, mundo y modelo.
    \item Se \concepto{representa} con distribuciones conjuntas y creencias.
    \item Se \concepto{actualiza} con Bayes ante nueva evidencia.
    \item Se \concepto{decide} comparando costo esperado de las acciones.
    \item Se \concepto{secuencializa} cuando las acciones cambian el futuro.
  \end{itemize}
\end{frame}

% 41
\begin{frame}{Errores que deben evitarse (repaso)}
  \small
  \concepto{Repaso:} esta lista integra los errores de razonamiento (diapositivas 9--10) y de decisión (diapositiva 28) ya estudiados.
  \begin{itemize}
    \item Confundir estado real, percepción y creencia.
    \item Invertir condicionales o ignorar la tasa base.
    \item Contar dos veces evidencia dependiente.
    \item Presentar una creencia como certeza.
    \item Valorar los errores de decisión como equivalentes.
    \item Decidir sin análisis de sensibilidad sobre costos y probabilidades.
    \item Aplicar RL a un problema que no es secuencial o donde explorar daña.
  \end{itemize}
  \vfill
  \alert{Una respuesta correcta pertenece a un modelo y un objetivo declarados; la responsabilidad es examinar ambos.}
\end{frame}

% 42
\begin{frame}{Cierre y próxima clase}
  \small
  Hoy construimos el puente entre el agente que percibe y el agente que decide.
  \begin{itemize}
    \item La incertidumbre es parte del problema, no un ruido a ignorar.
    \item La probabilidad representa creencias; la decisión agrega consecuencias.
    \item La secuencialidad prepara el terreno para el aprendizaje por refuerzo.
  \end{itemize}
  \vfill
  \begin{block}{Próxima clase}
    Actividad práctica integradora (robot y vehículo autónomo) y diseño experimental de aprendizaje automático: supervisado, no supervisado y evaluación fuera de muestra.
  \end{block}
\end{frame}

\end{document}
